\chapter{Jesus Cristo padeceu sob Pôncio Pilatos, foi crucificado, morto e sepultado}
\chaptermark{Jesus Cristo foi morto}

O quarto artigo da Profissão de Fé apresenta o núcleo do Mistério Pascal: a Paixão, a Morte e o Sepultamento de Jesus Cristo. ``Jesus Cristo padeceu sob Pôncio Pilatos, foi crucificado, morto e sepultado'' não descreve apenas um drama histórico, mas o cumprimento do plano salvífico de Deus. O Mistério Pascal de Cristo está no centro da Boa Nova que a Igreja anuncia ao mundo. Deus realizou a salvação ``de uma vez por todas'' pela morte redentora do seu Filho. A Cruz, que parecia o maior fracasso, revela-se como o supremo ato de amor e de obediência do Filho ao Pai.

\section{Introdução (571--573)}

O Mistério Pascal ocupa o centro da pregação apostólica. Os Apóstolos anunciaram ao mundo que Jesus, rejeitado pelos chefes do povo de Israel, foi entregue aos gentios, flagelado e crucificado. Contudo, esta morte não foi um acidente da história, mas o cumprimento das Escrituras. O próprio Jesus, depois da Ressurreição, explicou aos discípulos que era necessário que o Cristo sofresse estas coisas e entrasse na sua glória.

A Igreja permanece fiel à interpretação que o próprio Jesus deu de toda a Escritura. Os sofrimentos de Jesus assumiram uma forma histórica concreta: foi rejeitado pelos anciãos, pelos sumos sacerdotes e pelos escribas, que o entregaram aos romanos para ser flagelado e crucificado. Esta rejeição não foi universal. O povo simples de Israel muitas vezes seguia Jesus com admiração, enquanto as autoridades religiosas o viam como uma ameaça.

A fé pode, portanto, examinar as circunstâncias históricas da morte de Jesus, transmitidas fielmente pelos Evangelhos e iluminadas por outras fontes históricas, para compreender melhor o significado da Redenção. Longe de ser um episódio marginal, a Paixão de Cristo constitui o coração do Evangelho e a fonte da salvação para toda a humanidade.

A Cruz revela simultaneamente a gravidade do pecado humano e a imensidão do amor misericordioso de Deus. Ao contemplar o Crucificado, o cristão aprende a profundidade da sua própria miséria e a altura do amor divino que não recuou diante do sacrifício supremo.

\section{Jesus e Israel (574--594)}

Desde o início do ministério público de Jesus, certos fariseus, partidários de Herodes, sacerdotes e escribas decidiram eliminá-lo. Os milagres de expulsão de demónios, o perdão dos pecados, as curas em dia de sábado, a interpretação original da Lei sobre a pureza e a familiaridade com publicanos e pecadores provocaram forte oposição. Alguns chegaram a acusá-lo de estar possuído por um demónio.

Muitos gestos e palavras de Jesus constituíram um ``sinal de contradição'', especialmente para as autoridades religiosas de Jerusalém. O Evangelho de João refere-se frequentemente a ``os judeus'' ao falar destas autoridades. No entanto, as relações de Jesus com os fariseus não foram exclusivamente polémicas. Alguns o avisaram do perigo, Jesus elogiou alguns deles e aceitou convites para jantar em suas casas.

Aos olhos de muitos em Israel, Jesus parecia agir contra instituições essenciais do Povo Eleito: a observância integral da Lei, a centralidade do Templo e a fé no único Deus. No entanto, Jesus não veio abolir a Lei, mas cumpri-la perfeitamente. Ele declarou solenemente: ``Não penseis que vim abolir a Lei ou os Profetas. Não vim abolir, mas cumprir''.

Jesus manifestou o mais profundo respeito pelo Templo de Jerusalém. Apresentado nele ainda criança, aos doze anos permaneceu no Templo para tratar dos assuntos do seu Pai. Frequentou-o regularmente nas festas judaicas. Expulsou os vendilhões com zelo pelo Pai, declarando: ``Não façais da casa de meu Pai uma casa de comércio''. 

No entanto, Jesus anunciou a destruição do Templo, profetizando que não ficaria pedra sobre pedra. Com isto indicava que o seu próprio Corpo seria o novo Templo, o lugar definitivo da presença de Deus entre os homens. A sua morte corporal prefigurava a destruição do Templo de pedra e o início de uma nova era na história da salvação.

Jesus escandalizou sobretudo ao perdoar os pecados e ao identificar a sua conduta misericordiosa com a atitude do próprio Deus. Ao afirmar ``Eu e o Pai somos um'', colocava as autoridades religiosas diante da decisão fundamental: ou aceitavam a sua identidade divina ou o condenavam como blasfemo. O Sinédrio considerou-o merecedor da morte por blasfémia.

\section{Jesus morreu crucificado (595--623)}

A morte de Jesus na cruz não foi um mero acidente político. Foi o cumprimento do desígnio salvífico de Deus. O Pai entregou o seu Filho por amor à humanidade. O Filho ofereceu-se livremente por amor ao Pai e por amor a nós. A Cruz é, portanto, o supremo ato de obediência e de amor redentor.

O processo contra Jesus envolveu tanto as autoridades judaicas como o poder romano representado por Pôncio Pilatos. Embora Pilatos o tenha declarado inocente por três vezes, cedeu à pressão da multidão manipulada. A condenação à morte por crucificação, pena reservada aos escravos e rebeldes, revela a profundidade da humilhação assumida por Cristo.

A Cruz manifesta o escândalo e a loucura para os homens, mas é poder e sabedoria de Deus para os que são chamados. O Crucificado atrai todos a Si. Do alto da Cruz, Jesus reina como Rei universal, não com o poder da força, mas com o poder do amor que se entrega.

Jesus tomou sobre si a maldição da Lei para nos redimir dela. A sua morte redime as transgressões cometidas sob a primeira aliança. Ele é o Cordeiro de Deus que tira o pecado do mundo. Na Cruz, o servo sofredor de Isaías realiza plenamente a sua missão: ``o Senhor fez recair sobre ele a iniquidade de todos nós''.

A responsabilidade pela morte de Jesus não pode ser atribuída indiscriminadamente a todos os judeus de então nem aos judeus de hoje. A Igreja afirma que todos os pecadores foram os verdadeiros autores da Paixão de Cristo. Cada vez que pecamos, de certo modo crucificamos novamente o Senhor. Esta verdade deve levar à conversão e não ao ódio contra qualquer povo.

A morte na Cruz revela o preço do nosso resgate. Deus não nos salvou com qualquer preço, mas ``comprou-nos'' com o sangue do seu próprio Filho. Este sacrifício de amor infinito é fonte inesgotável de misericórdia para todos os que se aproximam de Cristo com confiança. A Cruz é o trono de onde o Senhor derrama as graças da Redenção.

A contemplação da Cruz transforma o cristão. Diante do Crucificado, o ódio, o orgulho e o egoísmo perdem o sentido. Surge a gratidão, a compaixão e o desejo de reparação. A Cruz torna-se o critério de toda a vida cristã e o caminho da verdadeira glória.

\section{Jesus Cristo foi sepultado (624--630)}

A morte de Jesus foi uma morte verdadeira. O seu corpo foi realmente separado da alma, embora a sua Pessoa divina permanecesse unida tanto à alma como ao corpo. Jesus experimentou plenamente a condição dos mortos, incluindo a descida à mansão dos mortos. O seu sepultamento constitui o elo real entre o seu estado passível antes da Páscoa e o seu estado glorioso depois da Ressurreição.

O sepultamento de Jesus cumpre as Escrituras. Isaías havia profetizado que o Servo seria sepultado com os ricos. José de Arimateia, membro respeitado do Sinédrio e discípulo secreto de Jesus, oferece o seu próprio túmulo novo. O corpo de Jesus é depositado com respeito e amor por algumas santas mulheres que o acompanhavam.

O corpo de Cristo no sepulcro não conheceu a corrupção. O poder divino preservou-o da decomposição, conforme o salmo: ``Não permitirás que o teu Santo veja a corrupção''. Este corpo santo, ainda unido à Pessoa divina, permanecia fonte de vida e santidade mesmo no silêncio do túmulo.

O sepultamento de Jesus tem grande significado sacramental. O Batismo significa descer ao túmulo com Cristo para morrer para o pecado e ressuscitar para uma vida nova. ``Fomos sepultados com ele pelo batismo na morte, para que, como Cristo foi ressuscitado dos mortos pela glória do Pai, assim também nós vivamos uma vida nova''.

O dia do sepultamento de Jesus é o grande sábado de Deus. Cristo repousa no túmulo depois de ter consumado a obra da salvação. Este repouso santo prefigura o repouso definitivo que Deus reserva aos seus filhos na vida eterna. É o mistério do Sábado Santo, dia de silêncio, espera e esperança.

A fé no sepultamento de Cristo convida-nos a entrar no silêncio e na confiança. Tal como o corpo de Jesus repousou no túmulo, o cristão é chamado a repousar no abandono confiante nas mãos do Pai, esperando a ressurreição. O sepulcro vazio tornar-se-á, no terceiro dia, o anúncio da vitória definitiva sobre a morte.